\section{Conclusion}
\label{sec:conclusion}

\subsection{Summary of findings}

This project trained and compared reinforcement learning agents on the ICU-Sepsis-v2 benchmark under two observation regimes and a risk-sensitive reformulation of the objective. The main results are summarised below.

\begin{itemize}
    \item \textbf{Tabular methods recover most but not all of the optimum.} Value Iteration achieves 78.8\% survival with analytical $J^* = 0.785$ under the fixed cohort difficulty. Q-Learning and SARSA reach 76.1\% and 74.3\% pooled survival across three seeds, with analytical optimality ratios of $0.877 \pm 0.006$ and $0.877 \pm 0.003$ respectively. The two algorithms are at parity to within seed noise. The residual 12 percentage points to the VI optimum is the joint cost of the constant step size and the $\epsilon$ floor; closing it requires schedule changes, not more episodes.
    \item \textbf{The exploration regime sweep inverts the textbook ranking.} At 50{,}000 episodes pure exploit reaches $J/J^* = 0.917$ by converging to the value of the always-no-treatment policy, which is unusually competitive in this MDP because action 0 has zero intensity penalty and is optimal for many low-severity states. The decay-to-0.05 schedule used in the main runs reaches only 0.829 at this budget and 0.877 at 200{,}000 episodes. Exploration design pays off slowly here.
    \item \textbf{Linear function approximation collapses to a single action.} Linear-Q converges across all three seeds to a 100\% always-(V=none, F=none) policy at 66.6\% survival, regardless of weight initialisation. The collapse is representational, not tie-breaking: the linear functional form cannot express the multi-modal optimal policy, and the always-action-0 fixed point is so competitive that the TD updates have weak incentive to escape it.
    \item \textbf{DQN learns a non-degenerate but underperforming policy.} 150{,}000-step training leaves the learning curve oscillating between 0.54 and 0.64 without convergence. The deployed best-snapshot policy reaches 67.4\% survival and shows clinically sensible SOFA-responsive vasopressor escalation, qualitatively matching VI's pattern. The deep approximator buys behaviour where the linear one cannot.
    \item \textbf{The wrappers cost the optimal policy almost nothing, but model-free policies lose more.} VI deployed on the clinical environment achieves 78.5\% survival, only 0.3 pp below clean VI. Q-Learning and SARSA tabular policies deployed via the same oracle protocol lose 4.5 and 3.9 pp respectively, an order of magnitude more. The 11.4 pp Configuration B gap to clean VI decomposes into 0.3 pp wrapper cost, 6.2 pp algorithm-quality gap, and 4.9 pp observation-regime cost.
    \item \textbf{Risk-sensitive shaping improves the centre, not the tail, with a sweet spot at $\beta = 0.5$.} Five-seed sweeps show survival rising from 72.8\% at $\beta = 0$ to 76.2\% at $\beta = 2$ ($\pm 0.5$--$1.1$ pp). Mean return peaks at $\beta = 0.5$ at 0.663 rather than rising monotonically. The 5\%-quantile of returns stays around $-0.12$ across all $\beta$ because the tail is dominated by acute events.
\end{itemize}

\subsection{Limitations}

Five limitations bound the generality of these results. \textbf{Simulator, not patients:} every number comes from ICU-Sepsis-v2, a fitted MDP. Policies that look good in simulation can still fail on patients whose dynamics differ from the fit, and no off-policy evaluation against held-out clinical trajectories was attempted. \textbf{Three seeds in Configuration A, five in the risk sweep:} cross-seed standard deviations are informative but underpowered for tight statistical claims. The Q-Learning vs SARSA parity claim would benefit from ten or more seeds. \textbf{DQN never converged:} 150{,}000 steps leave the curve oscillating; longer training, a smaller learning rate, soft target updates, or prioritised replay could plausibly close part of the gap to the oracle. The current DQN result is a lower bound on what deep RL could achieve. \textbf{Clinical realism of the wrappers:} the noise, missing-lab, and acute-event probabilities are reasonable proxies but not calibrated to real ICU error rates. \textbf{Action interpretability:} vasopressor and fluid levels are intuitive, but the policies do not explain why they chose a particular cell in a particular state; SHAP or input-gradient saliency on DQN, or coefficient-level inspection of Linear-Q, would close that gap.

\subsection{Potential improvements}

Two directions would extend this work without changing its scope. First, replace Linear-Q with a tile-coded or random-feature approximator: the failure mode is representational, and a richer (but still linear) feature space might break the always-action-0 fixed point cleanly. Second, retrain DQN with prioritised experience replay and a smaller learning rate to test whether the 150{,}000-step plateau is hyperparameter-bound rather than fundamental.
